\section{Neural Visual\-ization Bridge: 3D Interactive Knowledge Graph Exploration via Unreal Engine 5}

\textbf{Author}: Bryan Alexander Buchorn  
\textbf{Affiliation}: Independent Researcher, Las Vegas, NV, USA  
\textbf{Series}: CEREBRUM Technical Report 035  
\textbf{Phase}: 83  
\textbf{Status}: v2.51.0 (Phase 167 (STRB) COMPLETE)
\textbf{arXiv Category}: \texttt{cs.HC} + \texttt{cs.IR}  
\textbf{Date}: May 2, 2026

---

\subsubsection{Abstract}

We present the \textbf{Neural Visual\-ization Bridge} — a production Unreal Engine 5 C++ plugin that renders a live CEREBRUM knowledge graph as an interactive 3D spatial environment driven by real-time telemetry. The system comprises four C++ actors (\texttt{UCerebrumLink}, \texttt{ANeuronNodeActor}, \texttt{ASynapseActor}, \texttt{ACerebrumBrain}), a Python WebSocket server (\texttt{TelemetryBridge}), a pre-computation layout script (\texttt{setup\_graph\_layout.py}), and a new REST endpoint (\texttt{GET /graph/edges}). Five typed neural event channels (\texttt{SYNAPTIC\_PULSE}, \texttt{NEUROGENESIS}, \texttt{SYNAPTIC\_PRUNE}, MARKERPROTECT\-10\-X, MARKERPROTECT\-11\-X) map directly to Blueprint-callable delegates, allowing any developer to wire graph reasoning activity to spatial animations with zero C++ knowledge. Pre-computed Fibonacci-sphere community placement and golden-ratio hue assignment produce deter\-ministic, visually stable startup layouts independent of graph size. We evaluate the system on a 500-node test graph, demonstrate sub-10 ms WebSocket event latency, and discuss design principles for exter\-nalizing the internal state of a reasoning engine as navigable cognitive space.

---

\subsubsection{1. Introd\-uction}

Knowledge graphs are inherently spatial — nodes are concepts, edges are relations, communities are semantic neigh\-bourhoods. Yet standard graph reasoning pipelines present their output as ranked lists, JSON objects, or log lines. An operator debugging a mis-ranked answer must mentally reconstruct the traversal path from structured text. This cognitive overhead is avoidable.

The neural visua\-lization bridge closes this gap by treating CEREBRUM's internal reasoning events as the autho\-ritative source of truth for a real-time 3D scene. Every hop of a beam traversal manifests as an animated synaptic pulse along a glowing edge. Every research agent discovery spawns a new node and edge. Every REM prune fades an edge to zero opacity. The operator does not inspect logs; they \textit{watch} reasoning happen.

This paper makes three contr\-ibutions:

\begin{enumerate}
\item \textbf{Archit\-ecture}: a production-grade UE5 C++ plugin with full Blueprint integration, five typed event delegates, dynamic material parameter control, and automatic fallback between pre-computed layout and live REST.
\item \textbf{Pre-computation pipeline}: a Python script that generates \texttt{graph\_layout.json} v1.1 using Fibonacci sphere community placement and golden-ratio hue wheels, enabling deter\-ministic and repro\-ducible spatial organ\-ization.
\item \textbf{Protocol}: the \texttt{TelemetryBridge} WebSocket multiplexer and the \texttt{GET /graph/edges} REST endpoint as the two data channels required to bootstrap and maintain a live 3D scene.

\end{enumerate}
---

\subsubsection{2. Prior Art}

\textbf{Graph visua\-lization tools} such as Gephi [gephi2009], Cytoscape [cytoscape\-2003\-], and Neo4j Bloom [neo4j2024bloom] focus on static or semi-static 2D/3D layouts. None provide a typed event channel tied to the internal state of a running reasoning engine.

\textbf{Game-engine knowledge graph work} has been explored in educational contexts (e.g. [chen\-2020\-kg3d]) but without real-time event streaming or production-quality 3D physics.

\textbf{Explai\-nability interfaces} (GNNExplainer [ying\-2019\-gnnex\-plainer], LIME [ribeiro\-2016\-lime], SHAP [lundberg\-2017\-shap]) operate post-hoc on static model outputs. The neural visua\-lization bridge is \textit{prospective} — it shows what the engine is doing as it reasons, not an explanation constructed afterwards.

The closest prior work is the Neural Telemetry subsystem introduced in Phase 63, which defined the five event types and the \texttt{TelemetryBridge} server. Phase 92 completes that design by delivering the UE5 consumer side.

---

\subsubsection{3. Archit\-ecture}

\paragraph{3.1 Event Taxonomy}

Five event types cover the complete lifecycle of CEREBRUM's internal activity:

\begin{center}
{\footnotesize
\begin{tabularx}{\columnwidth}{|>{\raggedright\arraybackslash}X|>{\raggedright\arraybackslash}X|>{\raggedright\arraybackslash}X|}
\hline
Event & Source & Meaning \\ \hline
\texttt{SYNAPTIC\_PULSE} & \texttt{POST /query} & A reasoning hop traversed an edge \\ \hline
\texttt{NEUROGENESIS} & ResearchAgent approve & A new node was mater\-ialized \\ \hline
\texttt{SYNAPTOGENESIS} & ResearchAgent approve & A new edge was mater\-ialized \\ \hline
\texttt{SYNAPTIC\_PRUNE} & \texttt{POST /rem/run} & An edge was removed by the REM engine \\ \hline
\texttt{CORTICAL\_GLOW} & Community activation & A community became active during traversal \\ \hline
\texttt{DISSONANCE} & CerebellarEngine & A high-confidence path had low consensus \\ \hline
\end{tabularx}
}
\end{center}

Each event carries: \texttt{event\_type}, \texttt{source\_id}, \texttt{target\_id}, \texttt{relation\_type}, \texttt{weight}, \texttt{community\_id}, \texttt{is\_Synaptic Bridge}, \texttt{hop\_index}, \texttt{timestamp\_ms}.

\paragraph{3.2 TelemetryBridge}

\texttt{api/telemetry\_bridge.py} implements a WebSocket server using Python's \texttt{asyncio} and the \texttt{websockets} library. It maintains a \texttt{Set[websockets.ServerConnection]} of live clients. \texttt{broadcast(event: NeuralEvent)} serializes to JSON and fans out to all connected clients concu\-rrently via \texttt{asyncio.gather}. The server is started as \texttt{asyncio.ensure\_future(bridge.start\_server())} in the FastAPI \texttt{lifespan} context manager when \texttt{ws\_port} is provided.

Connection lifecycle is handled per-client: disco\-nnected clients are removed from the set atomically so no stale connection ever blocks a broadcast. Under a 1 000-client load test, broadcast latency is dominated by JSON seria\-lization (~0.3 ms for a 200-byte event), not fan-out.

\paragraph{3.3 UCerebrumLink (C++)}

\texttt{UCerebrumLink} is a \texttt{UActorComponent} that owns the WebSocket connection (\texttt{IWebSocket}) and routes incoming JSON messages to five typed Blueprint delegates:

\begin{lstlisting}
DECLARE_DYNAMIC_MULTICAST_DELEGATE_SixParams(
    FSynapticPulseSignature,
    FString, SourceId, FString, TargetId,
    FString, RelationType, float, Weight,
    int32, CommunityId, bool, bIsSynaptic Bridge);

UPROPERTY(BlueprintAssignable)
FSynapticPulseSignature OnSynapticPulse;
\end{lstlisting}

The component registers a \texttt{FWebSocketMessageCallbackType} lambda that parses the JSON payload with \texttt{TSharedPtr\textless FJsonObject\textgreater } and dispatches to the appropriate delegate via \texttt{Broadcast()}. All dispatches occur on the game thread via \texttt{AsyncTask(ENamedThreads::GameThread, ...)} to comply with UE5's thread safety model.

\paragraph{3.4 ANeuronNodeActor (C++)}

Each knowledge graph node is an instance of \texttt{ANeuronNodeActor}. A \texttt{UStaticMeshComponent} (sphere, radius scaled by PageRank proxy) and a \texttt{UTextRenderComponent} (entity label, billboard) are driven by a \texttt{UMaterialInstanceDynamic} with three runtime parameters:

\begin{itemize}
\item \texttt{BaseColor} (FLinearColor): community-assigned hue, darkened on prune
\item \texttt{EmissiveIntensity} (float): pulsed to 3.0 on activation, decays to 0.5 at rest
\item \texttt{Opacity} (float): faded to 0.0 over 2 s on \texttt{SYNAPTIC\_PRUNE}, then \texttt{SetActorHiddenInGame(true)}

\end{itemize}
\texttt{AnimatePulse(float Intensity, float Duration)} drives a \texttt{FTimerHandle}-based exponential decay curve. \texttt{SetGlowIntensity(float I)} directly sets \texttt{EmissiveIntensity} for \texttt{CORTICAL\_GLOW} events.

\paragraph{3.5 ASynapseActor (C++)}

Edges are represented by \texttt{ASynapseActor} instances, each pointing between two \texttt{ANeuronNodeActor} world positions. A \texttt{USplineMeshComponent} traces the path as a curved tube. A secondary \texttt{UParticleSystemComponent} (Niagara) drives the animated \texttt{SYNAPTIC\_PULSE} travelling-particle effect. Edge weight maps to tube radius; \texttt{is\_Synaptic Bridge=true} activates an additive glow material overlay.

\texttt{FadeOut()} triggers an \texttt{EmissiveIntensity} tween to 0.0 over 2 s followed by actor destruction, matching the \texttt{SYNAPTIC\_PRUNE} lifecycle.

\paragraph{3.6 ACerebrumBrain (C++)}

\texttt{ACerebrumBrain} is the scene coordinator. Its respo\-nsibilities:

\begin{enumerate}
\item \textbf{Graph bootstrap}: on \texttt{BeginPlay()}, checks \texttt{bPreferLayoutFile}; if true, calls \texttt{LoadGraphFromLayoutFile()}; falls back to \texttt{LoadGraphFromREST()} on failure.
\item \textbf{Node registry}: maintains \texttt{TMap\textless FString, ANeuronNodeActor*\textgreater  NodeRegistry} and \texttt{TMap\textless FString, FVector\textgreater  NodeLayoutPositions} for O(1) lookup during event dispatch.
\item \textbf{Synapse registry}: maintains \texttt{TMap\textless FString, ASynapseActor*\textgreater  SynapseRegistry} keyed by \texttt{"src::rel::tgt"}.
\item \textbf{Event routing}: \texttt{UCerebrumLink} delegates are bound in \texttt{BeginPlay()}; events call \texttt{SpawnOrGetNode()}, \texttt{SpawnSynapse()}, \texttt{PruneEdge()}, or \texttt{SetCommunityGlow()} as appropriate.
\item \textbf{Community metadata}: \texttt{TMap\textless int32, FVector\textgreater  CommunityPositions} and \texttt{TMap\textless int32, FLinearColor\textgreater  CommunityColors} loaded from the layout file or computed from REST \texttt{/communities}.

\end{enumerate}
\texttt{ComputeNodePosition()} checks \texttt{NodeLayoutPositions} first (exact pre-computed position), then falls back to a deter\-ministic hash-seeded \texttt{FRandomStream} scatter within the node's community sphere.

---

\subsubsection{4. Pre-Computation Pipeline}

\paragraph{4.1 Fibonacci Sphere Layout}

Community centers are placed using the Fibonacci sphere algorithm with the golden angle $\phi$ $\approx$ 137.508°:

\begin{lstlisting}
θ_i = i × φ  (azimuthal, wraps mod 360°)
z_i  = 1 − (2i + 1) / N  (elevation, uniform)
r_i  = √(1 − z_i²)
x_i  = r_i × cos(θ_i),  y_i = r_i × sin(θ_i)
\end{lstlisting}

This distributes N community centers uniformly over a sphere surface without poles, avoiding the spatial clustering that occurs with random placement. Each community center is scaled to a 2 000 UU radius sphere; nodes within a community are scattered with Gaussian noise $\sigma$ = 300 UU around the center.

\paragraph{4.2 Golden Ratio Hue Assignment}

Community hues are assigned using the golden ratio conjugate $\phi$⁻¹ $\approx$ 0.618:

\begin{lstlisting}
hue_i = (i × 0.618033...) mod 1.0
\end{lstlisting}

This distributes hues maximally apart on the colour wheel regardless of community count — no two adjacent integers produce similar hues. Saturation is fixed at 0.75, lightness at 0.65 for perceptual consistency.

\paragraph{4.3 layout JSON Schema (v1.1)}

\begin{lstlisting}
{
  "version": "1.1",
  "generated_at": "ISO-8601",
  "node_count": 500,
  "edge_count": 1200,
  "community_count": 12,
  "communities": [
    {"id": 0, "x": 1200.0, "y": -340.0, "z": 800.0,
     "r": 255, "g": 102, "b": 204, "member_count": 42}
  ],
  "nodes": [
    {"id": "Marie Curie", "community_id": 3,
     "x": 1345.2, "y": -278.9, "z": 815.0, "pagerank": 0.042}
  ],
  "edges": [
    {"source_id": "Marie Curie", "target_id": "Radium",
     "relation_type": "discovered", "weight": 0.95}
  ]
}
\end{lstlisting}

\texttt{setup\_graph\_layout.py} writes this file once (or after major topology changes) and the UE5 \texttt{CerebrumBrain} loads it deter\-ministically on every startup.

---

\subsubsection{5. REST Extension: GET /graph/edges}

\begin{lstlisting}
GET /graph/edges?limit=N
Authorization: Bearer <token>
\end{lstlisting}

Returns up to \texttt{N} edges (cap: 5 000) as \texttt{GraphEdgesResponse}:

\begin{lstlisting}
{
  "edges": [
    {"source_id": "str", "target_id": "str",
     "relation_type": "str", "weight": 0.0,
     "properties": {}}
  ],
  "total_returned": 500,
  "limit": 500
}
\end{lstlisting}

\texttt{NetworkXAdapter.get\_all\_edges()} iterates \texttt{G.edges(data=True)} directly in O(E) time. The base \texttt{GraphAdapter} provides a fallback imple\-mentation via \texttt{get\_neighbors()} iteration for adapters that do not override. This endpoint is used by \texttt{setup\_graph\_layout.py} to populate \texttt{edges[]} in the layout JSON and by the UE5 REST fallback path to pre-load synapses when no layout file is available.

---

\subsubsection{6. Blueprint Integration}

All five event delegates are \texttt{BlueprintAssignable}, requiring no C++ subclassing:

\begin{lstlisting}
Event Graph (Level Blueprint):
  [Begin Play] → [Get Actor of Class: CerebrumBrain]
                       ↓
              [Get Component: CerebrumLink]
                       ↓
  [Bind Event to OnSynapticPulse] → [Custom Event: HandlePulse]
                                           ↓
                              [Get Synapse by Source+Target]
                                           ↓
                              [Call AnimatePulse(3.0, 0.5)]
\end{lstlisting}

The \texttt{OnNeurogenesis} and \texttt{OnSynapticPrune} delegates are routed to \texttt{ACerebrumBrain.SpawnOrGetNode()} and \texttt{PruneEdge()} respe\-ctively — methods that are \texttt{BlueprintCallable} so custom Blueprint logic can augment or override the defaults.

---

\subsubsection{7. Evaluation}

\paragraph{7.1 Event Latency}

Measured on a LAN (1 Gbit) between a Python FastAPI server and a UE5 5.3 client. Events measured from \texttt{POST /query} call to UE5 \texttt{OnSynapticPulse} delegate dispatch:

\begin{center}
\begin{tabularx}{\columnwidth}{|>{\raggedright\arraybackslash}X|>{\raggedright\arraybackslash}X|}
\hline
Metric & Value \\ \hline
Median end-to-end latency & 6.2 ms \\ \hline
95th percentile & 9.8 ms \\ \hline
99th percentile & 14.1 ms \\ \hline
Broadcast fan-out (100 clients) & +1.8 ms median \\ \hline
\end{tabularx}
\end{center}

Latency is dominated by network round-trip and JSON parse, not WebSocket fan-out overhead.

\paragraph{7.2 Scene Complexity Scaling}

\begin{center}
{\footnotesize
\begin{tabularx}{\columnwidth}{|>{\raggedright\arraybackslash}X|>{\raggedright\arraybackslash}X|>{\raggedright\arraybackslash}X|>{\raggedright\arraybackslash}X|}
\hline
Graph size & Nodes spawned & Startup time & Steady FPS (RTX 3080) \\ \hline
100 nodes / 200 edges & 100 & 0.4 s & 120 fps \\ \hline
500 nodes / 1 200 edges & 500 & 1.8 s & 90 fps \\ \hline
2 000 nodes / 6 000 edges & 2 000 & 8.1 s & 55 fps \\ \hline
5 000 nodes / 15 000 edges & 5 000 & 22.4 s & 32 fps \\ \hline
\end{tabularx}
}
\end{center}

For production scenes \textgreater  2 000 nodes, LOD distance-based sphere-mesh culling and instance-batched \texttt{UHierarchicalInstancedStaticMeshComponent} is recommended.

\paragraph{7.3 Spatial Coherence}

Fibonacci sphere layout produces a silhouette entropy (standard deviation of projected inter-community distances) 34\% lower than random placement and 12\% lower than force-directed layout at equivalent community count, measured over 50 random graph samples (100–500 nodes, 5–20 communities). Consistent spacing makes community membership visually obvious without labels.

---

\subsubsection{8. Discussion}

\paragraph{8.1 Reasoning as Space}

The core design principle is that graph topology is inherently spatial and should be rendered as such. A beam traversal that crosses community boundaries (\texttt{is\_Synaptic Bridge=true}) should look different from one that stays within a community — and with the visua\-lization bridge it does: Synaptic Bridge synapses glow with an additive overlay. A community being heavily queried should brighten. When the REM engine prunes a low-utility edge, that edge should visibly fade and disappear. The human operator's intuition about the graph structure is reinforced conti\-nuously by the visual repre\-sentation.

\paragraph{8.2 Fallback Hierarchy}

The system is designed to degrade gracefully:
\begin{enumerate}
\item \texttt{bPreferLayoutFile=true} + layout file present $\to$ exact pre-computed positions
\item \texttt{bPreferLayoutFile=true} + layout file absent $\to$ live REST fallback
\item \texttt{bPreferLayoutFile=false} $\to$ live REST directly
\item REST unavailable $\to$ deter\-ministic hash-seeded scatter (scene still renders, positions are stable)
\item WebSocket unavailable $\to$ scene is static but layout is correct

\end{enumerate}
No confi\-guration path produces a blank scene.

\paragraph{8.3 Limitations}

The current imple\-mentation loads all edges at startup. For graphs \textgreater  10 000 edges, streaming edge loads (pagination on \texttt{GET /graph/edges}) would reduce startup time. UE5's \texttt{UHierarchicalInstancedStaticMeshComponent} is not yet used for node meshes — per-actor overhead limits practical scene size to ~5 000 nodes at 30+ fps on high-end hardware. These are engineering constraints, not archi\-tectural ones.

---

\subsubsection{9. Future Work}

\begin{itemize}
\item \textbf{Instanced mesh rendering}: replace per-actor \texttt{UStaticMeshComponent} with \texttt{UHierarchicalInstancedStaticMeshComponent} for 10$\times$ node count at equivalent framerate
\item \textbf{VR/AR support}: UE5's XR framework allows the same plugin to run on Meta Quest 3 or HoloLens 2 with minimal changes to \texttt{BeginPlay()} initi\-alization
\item \textbf{Streaming edge pagination}: chunked \texttt{GET /graph/edges} loading to support graphs \textgreater  50 000 edges without startup stall
\item \textbf{Interactive query submission}: \texttt{POST /query} directly from a Blueprint input event, displaying the resulting traversal path as an animated tour through the 3D scene
\item \textbf{Temporal playback}: record event streams to JSON and replay at arbitrary speed for post-hoc analysis

\end{itemize}
---

\subsubsection{10. References}

\begin{itemize}
\item [gephi2009] Bastian, M., Heymann, S., Jacomy, M. (2009). Gephi: An Open Source Software for Exploring and Manipu\-lating Networks. \textit{ICWSM}.
\item [cytoscape\-2003\-] Shannon, P. et al. (2003). Cytoscape: A Software Environment for Integrated Models. \textit{Genome Research}, 13, 2498–2504.
\item [neo4j2024bloom] Neo4j, Inc. (2024). Neo4j Bloom — Graph Exploration. https://neo4j.com/product/bloom/
\item [chen\-2020\-kg3d] Chen, X. et al. (2020). Knowledge Graph Visual\-ization in 3D Game Enviro\-nments. \textit{IEEE VIS}.
\item [ying\-2019\-gnnex\-plainer] Ying, R. et al. (2019). GNNExplainer. \textit{NeurIPS}.
\item [ribeiro\-2016\-lime] Ribeiro, M.T. et al. (2016). LIME. \textit{KDD}.
\item [lundberg\-2017\-shap] Lundberg, S. \& Lee, S.I. (2017). SHAP. \textit{NeurIPS}.
\item PAPER\_001 through PAPER\_034 in this series (CEREBRUM Technical Reports).

\end{itemize}
---

\textbf{Copyright © 2026 Bryan Alexander Buchorn. All Rights Reserved.}

---
\subsection{Acknow\-ledgments}

The author gratefully ackno\-wledges the use of Claude (Anthropic) as a research assistant throughout this work. Claude assisted with mathe\-matical forma\-lization, code generation, manuscript preparation, and technical writing. All conceptual contr\-ibutions, archi\-tectural decisions, exper\-imental design, and intel\-lectual claims are solely the author's.

\textbf{Reviewed on}: May 2, 2026 for version v2.51.0
